\section{Quasi-stationary variance of the binary Galton--Watson process}
\label{m:app:variance}

\cref{m:sec:condmean} establishes that the subcritical binary Galton--Watson
process, conditioned on survival, has limiting mean $1/A(p)$. The mean alone
does not establish quasi-stationarity: to show that the conditional
\emph{distribution} settles we need at least the second moment, and this
appendix computes it.

\subsection*{The second moment}

For a \rv{} $X$ the moment generating function is $\ex{\ee^{tX}}$, which in
terms of the \pgf{} is $\phi(\ee^t)$. Write $M_n(t) = \phi_n(\ee^t)$ for the
moment generating function of $Z_n$, so that moments are recovered by
\begin{equation}
\label{m:eq:mgfmoments}
\ex{Z_n^{\,\rho}} = \dda{^{\rho}}{t^{\rho}}M_n(t)\bigg|_{t=0} .
\end{equation}
For the binary offspring law \cref{m:eq:pgfbinary} the first-generation function
is
\begin{equation}
\label{m:eq:mgf1}
M_1(t) = \phi(\ee^t) = p\,\ee^{2t} + (1-p),
\end{equation}
whose second derivative at $t=0$ is $4p$, so $\ex{Z_1^2}=4p$. Reaching
generation $n$ requires a little more. Since $\phi_n$ is the $n$-fold
composition of $\phi$,
\begin{equation*}
M_n(t) = \phi\bigl(\phi(\dots\phi(\ee^t)\dots)\bigr),
\end{equation*}
and differentiating by the chain rule gives a product of $n+1$ factors,
\begin{equation}
\label{m:eq:dmgf1}
M_n'(t) = \phi'\bigl(\phi_{n-1}(\ee^t)\bigr)\cdot\phi'\bigl(\phi_{n-2}(\ee^t)\bigr)
\cdots\phi'\bigl(\phi_1(\ee^t)\bigr)\cdot\phi'(\ee^t)\cdot\ee^t .
\end{equation}
Now $\phi'(z) = 2pz$ for the binary law, so each factor collapses,
\begin{equation*}
\phi'\bigl(\phi_k(\ee^t)\bigr) = 2p\,\phi_k(\ee^t) = 2p\,M_k(t),
\qquad
\phi'(\ee^t) = 2p\,\ee^t ,
\end{equation*}
and \cref{m:eq:dmgf1} becomes
\begin{equation}
\label{m:eq:dmgf2}
M_n'(t) = (2p)^n\,M_{n-1}(t)\,M_{n-2}(t)\cdots M_1(t)\,\ee^{2t} .
\end{equation}
Before differentiating again, note that every factor on the right equals one at
$t=0$: $M_k(0)=1$ for all $k$, and $\ee^0=1$. That is what makes the second
derivative tractable. Differentiating the product gives a sum of $n$ terms, each
identical to \cref{m:eq:dmgf2} except that a single factor has been
differentiated, and at $t=0$ all the undifferentiated factors become one,
leaving
\begin{equation*}
M_n''(0) = (2p)^n\lrbs{M_{n-1}'(0) + M_{n-2}'(0) + \cdots + M_1'(0) + 2}.
\end{equation*}
But $M_k'(0)$ is the mean population at generation $k$, namely $(2p)^k$, so
\begin{equation*}
M_n''(0) = (2p)^n\lrbs{(2p)^{n-1} + \cdots + (2p) + 2}
= (2p)^n\lrbs{\sum_{i=0}^{n-1}(2p)^i + 1}.
\end{equation*}
Summing the finite geometric series and rearranging,
\begin{equation}
\label{m:eq:secondMoment}
\ex{Z_n^2} = (2p)^n\lrb{1+\frac{1}{1-2p}} - (2p)^{2n}\lrb{\frac{1}{1-2p}} .
\end{equation}
The minus sign is not optional: at $n=1$ the expression must return
$\ex{Z_1^2}=4p$, and it does.

\subsection*{Late-time variance}

From $\var{Z_n} = \ex{Z_n^2}-\ex{Z_n}^2$ and \cref{m:eq:secondMoment},
\begin{equation*}
\var{Z_n} = (2p)^n\lrb{1+\frac{1}{1-2p}} - (2p)^{2n}\lrb{\frac{1}{1-2p}} - (2p)^{2n},
\end{equation*}
and for $2p<1$ the terms in $(2p)^{2n}$ are negligible beside $(2p)^n$ at large
$n$, so $\var{Z_n}\to(2p)^n\bigl(1+1/(1-2p)\bigr)$. It is tempting to read
stationarity of the conditional variance straight off that, but the step is not
available: the conditional variance is not the unconditional one divided by
anything simple. Start instead from the conditional moments, using
\cref{m:eq:lateS} in each:
\begin{align*}
\ex{Z_n^2\mid Z_n>0} &= \frac{\ex{Z_n^2}}{S_n}
\;\xrightarrow{n\to\infty}\;
\frac{(2p)^n\lrb{1+\frac{1}{1-2p}}}{A(p)(2p)^n}
= \frac{1}{A(p)}\lrb{1+\frac{1}{1-2p}},\\[4pt]
\ex{Z_n\mid Z_n>0} &= \frac{\ex{Z_n}}{S_n}
\;\xrightarrow{n\to\infty}\;
\frac{1}{A(p)} .
\end{align*}
Both limits are finite, so the conditional variance tends to a constant:
\begin{equation}
\label{m:eq:condvar}
\var{Z_n\mid Z_n>0} \;\longrightarrow\;
\frac{1}{A(p)}\lrb{1+\frac{1}{1-2p}} - \frac{1}{A(p)^2} .
\end{equation}
The first two moments of the conditional population therefore both stabilise,
and the Yaglom theorem quoted in \cref{m:sec:condmean} supplies the full
quasi-stationary distribution whose first moments these are. Note that
\cref{m:eq:condvar}, like \cref{m:eq:qsmean}, is expressed entirely in terms of
$p$ and the single unknown constant $A(p)$, whose analysis is the subject of
\ChA.

\FloatBarrier
